\section{A faithful family and theory extension}\label{sec:family}

Let $\mathcal M_T$ be the set of all structures for the language of
arithmetic with underlying set $\N$ that satisfy $T$. Their arithmetic
operations need not be the standard ones, and the interpretation of a
numeral $\bar n$ need not be the carrier element $n$. Completeness and
the downward L\"owenheim--Skolem theorem imply
$\mathcal M_T\ne\varnothing$: every model of $T$ is infinite, and a
countably infinite model can be transported to the carrier $\N$.

Define the product category
\[
 \mathcal D_T=\prod_{M\in\mathcal M_T}\Pfn.
\]
An object is a family of sets $(A_M)_M$; a morphism is a family of partial
functions with the indicated sources and targets. Composition and the
restriction, range, restriction-product, and coproduct operations are coordinatewise.
The functors of Proposition~\ref{prop:model} combine to give
\[
 F_T:S_T\longrightarrow\mathcal D_T,\qquad
 F_T(X)=(X_M)_M,\quad F_T([e]_T)=(F_M[e]_T)_M.
\]

\begin{theorem}\label{thm:faithful}
The functor $F_T$ is faithful and preserves restrictions, zeros, binary
restriction products and coproducts, ranges, and finite joins of
restriction idempotents. The category
$\mathcal D_T$ is split and locally connected, satisfies weak choice,
and has inverse images preserving binary joins of restriction
idempotents. Moreover,
$F_T(1_\omega)\ne0$.
\end{theorem}
\begin{proof}
Preservation follows coordinatewise from Proposition~\ref{prop:model}.
If $[e]_T\ne[f]_T$, let $E$ denote the uniform graph-equality sentence
in \eqref{eq:uniform}. The theory $T\cup\{\neg E\}$ is consistent.
Completeness and downward L\"owenheim--Skolem give
$M\in\mathcal M_T$ in which the two graphs differ. Hence
$F_M([e]_T)\ne F_M([f]_T)$, so $F_T([e]_T)\ne F_T([f]_T)$.

A restriction idempotent in $\mathcal D_T$ is a family of partial
identities on subsets. Its splitting uses those subsets as the
coordinate sets. For an object $A=(A_M)_M$, take the partial identity
on a singleton $\{x\}\subseteq A_M$ at one coordinate $M$ and the zero
map elsewhere. This restriction idempotent is connected. These
restriction idempotents detect the definedness and value of every coordinate of a partial morphism. Hence
$\mathcal D_T$ is locally connected.

A partial monomorphism has injective coordinate functions: test its
defining property on objects with a singleton at one coordinate
and the empty set at every other coordinate.
The partial inverses of these injections form a section, proving weak
choice. Inverse images commute with unions coordinatewise. Finally,
$\mathcal M_T$ is nonempty and each $X_M$ is nonempty, so the
identity of $F_T(X)$ differs from zero.
\end{proof}

For consistent computably enumerable $T\supseteq\PA$, the source $S_T$
is not locally connected \cite[Corollary 4.4, pp.~657--658]{DPM}.
The connected restriction idempotents used in the preceding proof
belong to $\mathcal D_T$ and need not be images of source restriction
idempotents.

For an inclusion $T\subseteq T'$ of consistent theories, define
\[
 Q_{T,T'}:S_T\to S_{T'},\qquad
 Q_{T,T'}X=X,\quad Q_{T,T'}[e]_T=[e]_{T'}.
\]
Provability is monotone, so this is a well-defined covariant functor.
It preserves restrictions, zeros, ranges, binary restriction products
and coproducts, and satisfies $Q_{T,T'}(\tau_T)=\tau_{T'}$.
Since $\mathcal M_{T'}\subseteq\mathcal M_T$, restriction of families
defines a covariant projection
$\pi_{T,T'}:\mathcal D_T\to\mathcal D_{T'}$. Fixing the assignments
$X\mapsto X_M$ once for each arithmetic structure gives a strictly
commuting square:
\[
\begin{tikzcd}[column sep=large]
 S_T \arrow[r,"Q_{T,T'}"] \arrow[d,"F_T"'] &
 S_{T'} \arrow[d,"F_{T'}"]\\
 \mathcal D_T \arrow[r,"\pi_{T,T'}"'] & \mathcal D_{T'}.
\end{tikzcd}
\]
On a morphism $[e]_T$, both composites have, at every
$M\in\mathcal M_{T'}$, the partial function defined by interpreting
$G_e$ in $M$. For $T\subseteq T'\subseteq T''$, the quotient
functors satisfy $Q_{T',T''}Q_{T,T'}=Q_{T,T''}$, and the coordinate
projections satisfy
$\pi_{T',T''}\pi_{T,T'}=\pi_{T,T''}$.
The identity cases follow from the same definitions.

\subsection{Lifting in the syntactic category}

\begin{lemma}\label{lem:switch}
Let $c,c'$ be complementary restriction idempotents on $X$ in $S_T$.
There is a total morphism $s:X\to X+X$ with
\[
 sc=i_1c,\qquad sc'=i_0c'.
\]
The morphism $\chi_{X,X}^{-1}F_Ts$ is the unique section of
$\nabla_{F_TX}(F_Tc'+F_Tc)$, where
$\chi_{X,X}:F_TX+F_TX\to F_T(X+X)$ is the coproduct isomorphism.
\end{lemma}
\begin{proof}
Choose $\Sigma_1$ representatives $C(x)=\exists u\,R(x,u)$ and
$C'(x)=\exists v\,R'(x,v)$ with primitive-recursive witness tests.
The theory $T$ proves universally their disjointness and coverage.
For $j=0,1$, put $\iota_j(x)=2x+j$. The graph
$z=\iota_j(x)$ represents the coproduct injection
$i_j:X\to X+X$ on numerical codes.
Search increasing finite bounds for a witness to either predicate,
returning $\iota_1(x)$ on finding a $C$-witness and $\iota_0(x)$
on finding a $C'$-witness. Fix a priority for simultaneous discoveries.
Let $e_s$ be an index of this program. The finite-search verification
in $\PA$ proves that the program halts
exactly when $C(x)\lor C'(x)$. Under the disjointness assumption it also
proves the exact graph equation
\[
 G_{e_s}(x,z)\leftrightarrow
 (C(x)\land z=\iota_1(x))\lor(C'(x)\land z=\iota_0(x)).
\]
Thus $T$ proves this equation and universal halting. The equations for
the restrictions follow, giving the total morphism $s=[e_s]_T$ of $S_T$.

Interpret this proved graph identity in $M\models T$.
Under the canonical coproduct identification, $F_M(s)$ sends
$x$ to its second tagged copy when
$M\models C(\gamma_X^M(x))$ and to its first tagged copy when
$M\models C'(\gamma_X^M(x))$.
These conditions partition $X_M$, so $F_M(s)$ is the inverse of
$\nabla_{X_M}(F_M(c')+F_M(c))$.
Taking these inverses coordinatewise gives the unique section
asserted in the lemma.
\end{proof}

\begin{corollary}\label{cor:pair}
For every consistent $T\supseteq\PA$, the functor $F_T$ satisfies the
lifting hypothesis of Theorem~\ref{thm:local-lifting}.
Fix $U=\omega$ and a Turing morphism
$\tau:U\times U\rightharpoonup U$ in $S_T$.
Choose an isomorphism $b_0:U\to U+U$ in $S_{\PA}$ and put
$b=Q_{\PA,T}(b_0)$. Define
\[
 d=\tau\Delta_U,\qquad
 a_0=(bd)^*(1_U+0_U),\qquad a_1=(bd)^*(0_U+1_U).
\]
Then $a_0,a_1$ are inseparable in $S_T$. For every consistent
extension $T'\supseteq T$, their images under $Q_{T,T'}$ are the
identically defined inseparable pair in $S_{T'}$.
\end{corollary}
\begin{proof}
Theorem~\ref{thm:faithful} and Lemma~\ref{lem:switch} verify all the
hypotheses of Theorem~\ref{thm:local-lifting}.
The functor $Q_{T,T'}$ preserves the operations defining the pair.
It is the identity on objects and surjective on each hom-set.
Consequently it preserves a Turing morphism: a preimage of a
target morphism has a Turing factorization in $S_T$, whose image is
the required factorization in $S_{T'}$.
The same lifting argument therefore applies over $T'$.
\end{proof}

The morphism $F_T(\tau_T)$ is not a Turing morphism in $\mathcal D_T$.
Fix $M\in\mathcal M_T$. The partial functions
$x\mapsto F_M(\tau_T)(a,x)$, indexed by all $a\in|M|$, form a
countable family, whereas there are uncountably many partial
functions $|M|\rightharpoonup|M|$.
Choose one outside this family and use it at coordinate $M$,
with zero partial functions at the other coordinates.
The resulting morphism
$\mathbf1\times F_T\omega\rightharpoonup F_T\omega$,
where $\mathbf1$ is the family of singleton sets, cannot factor
through $F_T(\tau_T)$ with a total parameter
$\mathbf1\to F_T\omega$.

The target restriction lattices are Boolean, so every image of a disjoint source
pair has a separator. Lemma~\ref{lem:switch} lifts the sections
associated with complementary source restriction idempotents;
an arbitrary target separator need not be the image of such an
idempotent.

